% !TEX program = pdflatex
% Documentation fonction par fonction de model_test_prototype.ipynb
% Compilation : pdflatex model_test_prototype_documentation.tex (2 passes pour les références)

\documentclass[11pt,a4paper]{article}

\usepackage[utf8]{inputenc}
\usepackage[T1]{fontenc}
\usepackage[french]{babel}
\usepackage{lmodern}
\usepackage[margin=2.2cm]{geometry}
\usepackage{booktabs}
\usepackage{array}
\usepackage{longtable}
\usepackage{xcolor}
\usepackage[hidelinks]{hyperref}

% Désactive les espacements automatiques du français devant ; : ! ?
\shorthandoff*{;:!?}

\newcommand{\code}[1]{\texttt{\small #1}}

\title{Documentation de \code{model\_test\_prototype.ipynb}}
\author{}
\date{}

\begin{document}
\maketitle

%======================================================================
\section{Hypothèses métier}
\label{sec:hypotheses}

{\small
\begin{longtable}{@{}l p{2.45cm} p{5.9cm} p{4.5cm}@{}}
\toprule
\textbf{Id} & \textbf{Étape} & \textbf{Hypothèse métier} & \textbf{Si elle est fausse} \\
\midrule
\endfirsthead
\toprule
\textbf{Id} & \textbf{Étape} & \textbf{Hypothèse métier} & \textbf{Si elle est fausse} \\
\midrule
\endhead
\bottomrule
\endfoot

H1 & \url{extract_all_triggers} & Un \code{exec\_code} correspond à un seul
\code{trigger\_internal\_id} : \code{dropDuplicates(["exec\_code"])} n'a rien à départager. &
Le \code{trigger\_internal\_id} retenu est arbitraire, et les jointures avals portent sur le
mauvais produit déclencheur. \\
\addlinespace
H2 & \url{extract_all_triggers} & Un trigger n'est évaluable que si son produit déclencheur est un
nœud du graphe de co-vue (\code{trigger\_internal\_id in node2idx}). & Les triggers hors catalogue
actif sortent silencieusement de la population évaluée ; seul \code{skipped\_no\_node} le révèle. \\

\end{longtable}
}
\refstepcounter{table}\label{tab:hypotheses}
\begin{center}\footnotesize\itshape Tableau \thetable{} --- Hypothèses métier.\end{center}

%======================================================================
\clearpage
\section{Ce qui peut être exclu de l'évaluation}
\label{sec:exclusions}

{\small
\begin{longtable}{@{}l p{3.2cm} p{4.4cm} p{5.2cm}@{}}
\toprule
\textbf{Id} & \textbf{Étape} & \textbf{Mécanisme} & \textbf{Ce qui disparaît} \\
\midrule
\endfirsthead
\toprule
\textbf{Id} & \textbf{Étape} & \textbf{Mécanisme} & \textbf{Ce qui disparaît} \\
\midrule
\endhead
\bottomrule
\endfoot

\end{longtable}
}
\refstepcounter{table}\label{tab:exclusions}
\begin{center}\footnotesize\itshape Tableau \thetable{} --- Ce qui peut être exclu de l'évaluation.\end{center}

\end{document}
